% Required packages: graphicx, subcaption, booktabs.

\section{Experiments}

\subsection{Main Comparison}

\paragraph{Experimental setting.}
We compare preference tracing against three GPT-5-based online personalization baselines on
the same matched users. The CoT baseline conditions on up to five previous rounds of dialogue
and feedback as in-context examples, and learns to imitate the accepted answer. The RAG
baseline stores observed dialogues in a vector database and retrieves the top five examples
for each new user message. The dynamic cheatsheet baseline gives the model a freely editable
user-profile notebook, which can be updated or extended after each dialogue and selection.
Preference tracing instead maintains an explicit preference state and uses it for both
response adaptation and preference prediction.

\begin{figure*}[t]
    \centering
    \begin{subfigure}[t]{0.49\textwidth}
        \centering
        \includegraphics[width=\linewidth]{figures/main_comparison_placeholder.png}
        \caption{Online trajectories.}
        \label{fig:main-comparison-traj}
    \end{subfigure}
    \hfill
    \begin{subfigure}[t]{0.49\textwidth}
        \centering
        \small
        \resizebox{\linewidth}{!}{
        \begin{tabular}{lrrrrrrrrrrrr}
        \toprule
        Method & $N$ & Acc.$_{>20}$ & $\Delta$Acc. & GPT-RA$_{>20}$ & $\Delta$GPT-RA & Emb.-RA$_{>20}$ & $\Delta$Emb.-RA & Survey & Aspect & Plaus. & Overall & Sim. \\
        \midrule
        GPT-5 trace & 28 & 0.6134 & 0.4211 & 0.2267 & 0.9190 & -0.0010 & 0.0221 & 4.1429 & 3.8929 & 4.8571 & 4.1857 & 0.5366 \\
        CoT & 28 & 0.5481 & 0.0865 & 0.0425 & 0.6579 & 0.0222 & 0.0549 & -- & -- & -- & -- & -- \\
        RAG & 28 & 0.5808 & 0.1962 & 0.1248 & 0.8556 & 0.0363 & 0.1162 & -- & -- & -- & -- & -- \\
        Cheatsheet & 28 & 0.4857 & 0.0242 & -0.2464 & 0.2152 & -0.0136 & 0.0130 & -- & -- & -- & -- & -- \\
        \bottomrule
        \end{tabular}
        }
        \caption{Aggregate metrics.}
        \label{fig:main-comparison-table}
    \end{subfigure}
    \caption{
    Main comparison on the matched-user evaluation split. The left panel shows turn-wise
    prediction accuracy, relative GPT response alignment, relative embedding alignment, and
    profile alignment. The right panel reports per-user means after turn 20 and their change
    from turn 0; Acc. denotes preference-prediction accuracy, GPT-RA and Emb.-RA denote
    relative response-alignment scores, and profile columns report rubric-based profile
    alignment.
    }
    \label{fig:main-comparison}
\end{figure*}

\paragraph{Results.}
Figure~\ref{fig:main-comparison} shows that preference tracing improves more steadily than
the baselines and obtains the highest late-stage prediction accuracy. After turn 20,
preference tracing reaches 0.6134 accuracy, compared with 0.5808 for RAG, 0.5481 for CoT,
and 0.4857 for the dynamic cheatsheet baseline. Its relative GPT alignment is also strongest
in the late stage, while the cheatsheet baseline remains negative. CoT has comparatively high
response alignment during parts of the online process because raw previous responses are
available as demonstrations, but this does not produce a stable accuracy-improvement signal.
RAG is competitive on embedding alignment, suggesting that retrieval helps recover
surface-similar evidence, but it remains below preference tracing on choice prediction.
\emph{Placeholder: add the cross-benchmark comparison conclusion once the second benchmark
result is finalized.}

\subsection{Model Ablation}

\paragraph{Experimental setting.}
We next evaluate whether preference tracing depends on a particular strong proprietary
backbone. We run the same tracing procedure with GPT-5, kimi-k2.6, GLM-5.1, DeepSeek-V4,
Qwen3.5-397B-A17B, and Qwen3.5-9B. All variants use Gemini-3-flash as the LLM critic for
response-alignment and profile-alignment evaluation, so differences reflect the tracing
backbone rather than changes in the judge.

\begin{figure*}[t]
    \centering
    \begin{subfigure}[t]{0.49\textwidth}
        \centering
        \includegraphics[width=\linewidth]{figures/model_ablation_placeholder.png}
        \caption{Online trajectories.}
        \label{fig:model-ablation-traj}
    \end{subfigure}
    \hfill
    \begin{subfigure}[t]{0.49\textwidth}
        \centering
        \small
        \resizebox{\linewidth}{!}{
        \begin{tabular}{lrrrrrrrrrrrr}
        \toprule
        Model & $N$ & Acc.$_{>20}$ & $\Delta$Acc. & GPT-RA$_{>20}$ & $\Delta$GPT-RA & Emb.-RA$_{>20}$ & $\Delta$Emb.-RA & Survey & Aspect & Plaus. & Overall & Sim. \\
        \midrule
        GPT-5 & 28 & 0.6134 & 0.4211 & 0.2267 & 0.9190 & -0.0010 & 0.0221 & 4.1429 & 3.8929 & 4.8571 & 4.1857 & 0.5366 \\
        kimi-k2.6 & 28 & 0.4904 & 0.1827 & 0.0659 & 0.9505 & -0.0020 & 0.0592 & 3.0000 & 2.7500 & 4.5000 & 3.2000 & 0.5251 \\
        GLM-5.1 & 28 & 0.5485 & 0.1639 & 0.0621 & 0.2544 & 0.0010 & 0.0546 & 3.2500 & 3.1071 & 4.6429 & 3.4714 & 0.5685 \\
        DeepSeek-V4 & 28 & 0.6089 & 0.3012 & 0.2378 & 1.6609 & -0.0053 & 0.0464 & 3.4286 & 3.1786 & 4.5357 & 3.5500 & 0.5129 \\
        Qwen3.5-397B-A17B & 28 & 0.5398 & 0.2706 & -0.0110 & 0.7582 & -0.0076 & 0.0496 & 2.8929 & 2.8214 & 4.1429 & 3.1143 & 0.5317 \\
        Qwen3.5-9B & 28 & 0.5211 & 0.1750 & 0.0052 & 0.4475 & -0.0039 & 0.0473 & 3.1429 & 3.1071 & 4.3214 & 3.3643 & 0.5144 \\
        \bottomrule
        \end{tabular}
        }
        \caption{Aggregate metrics.}
        \label{fig:model-ablation-table}
    \end{subfigure}
    \caption{
    Model ablation under a fixed Gemini-3-flash critic. The left panel compares online
    trajectories across tracing backbones, and the right panel reports the same late-stage
    aggregate metrics as Figure~\ref{fig:main-comparison}.
    }
    \label{fig:model-ablation}
\end{figure*}

\paragraph{Results.}
Figure~\ref{fig:model-ablation} shows that preference tracing yields useful online gains
across recent open and closed models. All backbones improve in prediction accuracy relative
to turn 0, and most maintain near-positive or positive late-stage relative response
alignment. This suggests that the method decomposes personalization into capabilities that
modern models broadly possess: simulating basic user preference choices, integrating evidence
over interactions, and assigning stable comparative scores. The improvement therefore does
not appear to rely on exceptionally strong long-horizon reasoning from the backbone.

The ablation also reveals non-monotonic behavior with model size. Qwen3.5-9B is more stable
than Qwen3.5-397B-A17B in late-stage relative GPT alignment and profile overall score,
despite being much smaller. A plausible explanation is that scoring stability is especially
important in the tracing loop: a larger model with less consistent comparative judgments may
update the preference state less reliably. At the same time, GPT-5 remains strongest on
profile alignment, indicating that accurate profile extraction still benefits from stronger
domain-sensitive abstraction even when the online choice-improvement mechanism transfers
across smaller models.
